由例 \ref{no singularity} 已经可见，$\Sht^\loc_\mu$ 不是代数叠，因此 \S \ref{SS: cat of corr} 中的一般构造不能直接用于当前情形。为了在此情形定义范畴，我们需要用有限维逼近替代这些模空间。
这里只勾勒构造，细节参见 \cite[\S 5]{XZ}。构造的群论描述见注 \ref{R: grp theoretical Sht}。

回忆 $K$ 作用在卷积 Grassmann 簇上。给定支配余权序列 $\mmu=(\mu_1,\ldots,\mu_r)$，称整数 $m$ 为 $\mmu$-大，若它 $\geq \langle\sum_i\mu_i,\al_h\rangle$，其中 $\al_h$ 是 $G$ 的最高根。
容易看出，给定非负整数对 $(m,n)$（允许 $(m,n)=(\infty,\infty)$），若 $m-n$ 是 $\mmu$-大，则 $K$ 在 $\Gr_\mmu^{(n)}$ 上的作用经过 $K_{m}$ 的作用分解（参看 \cite[引理 3.1.7]{XZ}）。此时 fpqc 商叠
$\Hk^{\loc(m)}_\mmu:=[K_m\backslash \Gr_\mmu]$ 称为 $m$-限制局部 Hecke 叠。
由于
\[
[K_m\backslash\Gr_\mmu]\cong [K_m\backslash\Gr^{(n)}_\mmu/K_n]\xrightarrow{t_\leftone\times t_\rightone} \bfB K_n\times \bfB K_m,
\]
我们看到，在 $\Hk_{\mmu}^{\loc(m)}$ 之上，第一个映射 $t_\rightone$ 定义 $K_m$-主丛
\begin{equation*}\label{torsor2}
\Gr_\mmu \to [K_m\backslash\Gr_\mmu],
\end{equation*}
第二个映射 $t_\leftone$ 定义 $K_n$-主丛
\begin{equation*}\label{torsor1}
[K_m\backslash\Gr_{\mmu}^{(n)}]\to  [K_m\backslash\Gr_\mmu].
\end{equation*}
把这两个\emph{典范}主丛分别记作 $\calE_{\rightone}|_{D_{m}}$ 与 $\calE_{\leftone}|_{D_n}$。

对 $m\geq n$，记
$$\res^m_n: \bfB K_m\to \bfB K_n$$
为由商映射 $\pi_{m,n}:K_m\to K_n$ 诱导的映射。显然，当 $m\leq m'\leq m''$ 时
\begin{equation}
\label{E:ex and q compatibility}
\res_{m}^{m'} \circ \res_{m'}^{m''} = \res_{m}^{m''},
\end{equation}
这里允许 $m''=\infty$。
于是对 $m\leq m'$、$n\leq n'$（允许 $(m',n')=(\infty,\infty)$），若 $m'-n'$ 与 $m-n$ 都是 $\mmu$-大，则有
限制映射
$$\res^{m'}_m: \Hk^{\loc(m')}_{\mmu}\to \Hk^{\loc(m)}_{\mmu},$$ 使得下图交换且右方块是笛卡儿的。
\begin{equation*}
\label{E:restriction and torsor commutative}
\xymatrix{
\bfB K_{n'} \ar[d]^{\res^{n'}_n} &
\Hk_\mmu^{\loc(m')} \ar[l]_-{t_\leftone} \ar[r]^-{t_\rightone} \ar[d]^{\res^{m'}_{m}} &  \bfB K_{m'} \ar[d]^{\res^{m'}_{m}}
\\
\bfB K_n & \Hk_\mmu^{\loc(m)} \ar[l]_-{t_\leftone} \ar[r]^-{t_\rightone} &  \bfB K_m.
}
\end{equation*}
此外，这些限制映射满足自然的相容性条件 \eqref{E:ex and q compatibility}。

现在固定支配余权序列 $\mmu$，并选整数对 $(m,n)$ 使 $m-n$ 为 $\mmu$-大。
对完美 $k$-代数 $R$，$\Spec R$ 上奇点被 $\mmu$ 所界的 $(m,n)$-限制局部迭代 $G$-Shtuka 由以下数据组成：
\begin{itemize}
\item $ \Hk^{\loc(m)}_\mmu$ 的一个 $R$-点，以及
\item $\Spec R$ 上 $K_n$-主丛的同构
$$\psi:{^\sigma} (\mE_\leftone|_{D_{n}}) \simeq(\mE_\rightone|_{D_{m}})|_{D_{n}}.$$
\end{itemize}
把分类所有 $(m,n)$-限制局部 $G$-Shtuka 的预叠定义为 $\Sht^{\loc(m,n)}_\mmu$。换言之，$\Sht_\mmu^{\loc(m,n)}$ 定义为 $\Hk_\mmu^{\loc(m)}$ 与一个 Frobenius 图的纤维积
\begin{equation}
\label{E:mS restricted as a Cartesian product}
\xymatrix@C=90pt{
\Sht_\mmu^{\loc(m,n)} \ar[r]^{\varphi^{\loc(m,n)}} \ar[d]
&  \Hk_\mmu^{\loc(m)} \ar[d]^-{ t_\leftone\times \res^{m}_{n}\circ t_\rightone }
\\ \bfB K_n \ar[r]^{1\times \sigma} & \bfB K_n \times \bfB K_n,
}
\end{equation}
特别地，$\Sht_\mmu^{\loc(m,0)}=\Hk_\mmu^{\loc(m)}$，且 $\Sht_\mmu^{\loc(\infty,\infty)}=\Sht_\mmu^\loc$。
注意 $\varphi^{\loc(m,n)}$ 是完全光滑的态射，相对维数为 $n \dim G$。

设 $(m',n')$ 与 $(m,n)$ 是两对非负整数，满足 $m\leq m'$、$n\leq n'$，且 $m'-n'$ 与 $m-n$ 都是 $\mmu$-大（允许 $(m',n')=(\infty,\infty)$）。
则有\emph{限制态射}
\begin{equation}\label{trunc}
\res^{m',n'}_{m,n}: \Sht^{\loc(m',n')}_\mmu\longto \Sht_\mmu^{\loc(m,n)},
\end{equation}
它是相对维数 $(m-n)-(m'-n')$ 的完全光滑态射，与 \eqref{E:mS restricted as a Cartesian product} 相容，并满足
\begin{equation}
\label{true-to-trun}
\res^{m_2,n_2}_{m_1,n_1}\circ\res^{m_3,n_3}_{m_2,n_2}=\res^{m_3,n_3}_{m_1,n_1}.
\end{equation}

\begin{ex}
\label{Ex:mStau mn}
当 $\mu=0$，更一般地当 $\mu=\tau\in\xcoch(Z_G)$ 时，有
$$\Sht^{\loc(n,n)}_\tau\cong \bfB G(\mO/\varpi^n).$$
限制映射 $\res_{n,n}: \Sht_\tau^\loc\cong\bfB G(\mO)\to \Sht_\tau^{\loc(m,n)}\cong\bfB G(\mO/\varpi^n)$ 由自然映射 $G(\mO)\to G(\mO/\varpi^n)$ 诱导。
\end{ex}

\begin{rmk}
\label{Ex:Sht and GZip}
如例 \ref{loc Sht and pdiv} 所述，局部 Shtuka 的模空间可以看作带 $G$-结构的 $p$-可除群模空间的推广。自然期待：限制局部 Shtuka 的模空间可以看作带 $G$-结构的截断 Barsotti–Tate 群模空间的推广。结果表明两者确实密切相关，但关系比非限制情形略复杂。如 \cite[引理 5.3.6]{XZ} 所述，若 $\mu$ 极小且 $(m,n)=(2,1)$，则有自然完全光滑映射
\begin{equation}
\label{E:Sht-Zip}
\Sht_\mu^{\loc(2,1)}\to G\on{-Zip}^\pf_\mu.
\end{equation}
这里 $G\on{-Zip}^\pf_\mu$ 是 \cite{PWZ} 中定义的带 $G$-结构的 $F$-zip 模空间的完全化，它同胚于带 $G$-结构的 $1$-截断 Barsotti–Tate 群的模空间。

我们描述 $G=\GL_n$、$\mu=\omega_i$ 时的映射 \eqref{E:Sht-Zip}。此时 $G\on{-Zip}_{\omega_i}(R)$ 分类四元组 $(M, C^\bullet M, D_\bullet M, \phi_0,\phi_1)$，其中 $M$ 是秩 $n$ 的有限投射 $R$-模，$M=C^0M\supset C^1M\supset C^2M=0$ 是递减滤过，$D_{-1}M=0\subset D_0M\subset D_1M$ 是递增滤过，$\phi_0:\sigma^*\on{gr}_C^0M\cong \on{gr}_0^DM$ 是秩 $i$ 的有限投射 $R$-模之间的同构，$\phi_1: \sigma^*\on{gr}_C^1M\cong \on{gr}_1^DM$ 是秩 $n-i$ 的有限投射 $R$-模之间的同构。上述类型的 $F$-zip 与高度 $n$、维数 $i$ 的 $1$-截断 Barsotti–Tate 群之间的关系见 \cite[\S 9.3]{PWZ}。

注意给出 $\Sht_{\omega_i}^{\loc(2,1)}$ 的 $R$-点，就有一一对应的有限投射 $R$-模映射
$$\mE_\leftone|_{R}\xrightarrow{\beta} \mE_\rightone|_{R}\cong \sigma^*(\mE_\leftone|_{R}),$$ 其余核是秩 $i$ 的有限投射 $R$-模。
于是 \eqref{E:Sht-Zip} 把 $\Sht_{\omega_i}^{\loc(2,1)}$ 的 $R$-点送到 $G\on{-Zip}_{\omega_i}$ 的 $R$-点，后者由 $M=\mE_\rightone|_{R}$、$C^1M=\on{Im} \beta$、$D_0M=\ker(\sigma^*\beta)$ 表示。同构 $\phi_1: \sigma^*(\on{gr}_C^1M)=\sigma^*\on{Im}\beta\cong \on{gr}^D_1M=\frac{M}{\ker\sigma^*\beta}$ 是典范的。
典范同构 $\phi_0: \sigma^*(\on{gr}_C^0M)=\sigma^*\frac{M}{\on{Im}\beta}\cong \ker(\sigma^*\beta)=\on{gr}^0_DM$ 可由 \cite[(1.3.2)]{Zh1} 推出。
\end{rmk}

各种限制局部迭代 Shtuka 模空间之间存在 Hecke 对应，它们经限制态射 \eqref{trunc} 与对应 \eqref{E;Hk Shtuka} 相容。此外，这些 Hecke 对应可以复合，且复合满足结合律。我们将在命题 \ref{P:pushforward is local} 中总结这些结构，它使我们能够在下一小节定义范畴 $\on{P}^{\Hk}(\Sht^\loc_{\bar k})$。

引入下面的术语会很方便。
给定支配余权序列 $\mmu=(\mu_1,\ldots,\mu_t)$，若非负整数四元组 $(m_1,n_1,m_2,n_2)$ 满足
\begin{enumerate}
\item
$m_1-m_2=n_1-n_2$ 是 $\mu_t$-大（等价地，是 $\sigma(\mu_t)$-大），
\item
$m_2-n_1$ 是 $(\mu_1, \dots, \mu_{t-1})$-大，
\end{enumerate}
则称其为 $\mmu$-\emph{可接受}的。特别地，$m_1-n_1$ 是 $\mmu$-大。我们把 $(\infty,\infty,\infty,\infty)$ 看作对任意 $\mmu$ 都可接受。

下面的命题总结了 \cite[\S 5.3]{XZ} 的大部分内容。
为简便，我们只考虑奇点被单个余权所界的限制局部 Shtuka 的模空间。
